\section{Preliminaries}
\subsection{Gaussian surfels model}
Our PolGS selects \gs as our base framework due to its stronger geometry expression ability. 
According to \gs, we use a set of unstructured Gaussian kernels $\{G_i=\{\textbf{x}_i,\textbf{s}_i,\textbf{r}_i,o_i,C_i\}|i\in
 \mathcal{N}\}$ to represent the structure of 3DGS, where  $\textbf{x}_i\in\mathbb{R}^3$ denotes the center 
 position of each Gaussian kernel, $\textbf{s}_i=[s^x_i,s^y_i,0]^\top\in\mathbb{R}^3$ is the scaling factors of $x$ and 
 $y$ axes after flatting the 3D Gaussians~\cite{kerbl20233d}, $\textbf{r}_i\in\mathbb{R}^4$ is the rotation quaternion, 
 $o_i\in\mathbb{R}$ is the opacity, and $C_i\in\mathbb{R}^k$ represents the spherical harmonic coefficients of each Gaussian. 
 And Gaussian distribution can be defined by the covariance matrix $\Sigma$ of a 3D Gaussian as:
\begin{equation}
    G(\textbf{x};\textbf{x}_i,\Sigma_i)=\text{exp}(-\frac{1}{2}(\textbf{x}-\textbf{x}_i)^\top\Sigma_i^{-1}(\textbf{x}-\textbf{x}_i)),
\end{equation}
where $\Sigma_i$ can be represented as:
\begin{equation}
\begin{aligned}
	\label{equation:sigma}
    \Sigma_i &= \textbf{R}(\textbf{r}_i)\textbf{s}_i\textbf{s}_i^\top\textbf{R}(\textbf{r}_i)^\top\\
    &= \textbf{R}(\textbf{r}_i)\text{Diag}[(s^x_i)^2,(s^y_i)^2,0]\textbf{R}(\textbf{r}_i)^\top,
\end{aligned}
\end{equation}
where Diag[·] indicates a diagonal matrix and $\textbf{R}(\textbf{r}_i)$ is a $3 \times 3$
rotation matrix represented by $\textbf{r}_i$.
\vspace{-1em}
\paragraph{Gaussian splatting}
According to 3D GS~\cite{kerbl20233d}, novel view rendering process can be represented as: 
\begin{equation}
    C=\sum_{i=0}^nT_i\alpha_ic_i,
\label{eq:gs_render}
\end{equation}
where $T_i=\prod_{j=0}^{i-1}(1-\alpha_j)$ is the transmittance, $\alpha_i = G'(\textbf{u};\textbf{u}_i,\Sigma'_i)o_i$ 
is alpha-blending weight, which is the product of opacity and the Gaussian weight based on pixel $\textbf{u}$. 
In order to speed up the rendering process, the 3D Gaussian in \eref{equation:sigma} 
is re-parameterized in 2D ray space~\cite{zwicker2002ewa} as $G'$:
\begin{equation}
    G'(\textbf{u};\textbf{u}_i,\Sigma'_i)=\text{exp}(-\frac{1}{2}(\textbf{u}-\textbf{u}_i)^\top {\Sigma'_i}^{-1}(\textbf{u}-\textbf{u}_i)),
\end{equation}
where ${\Sigma'_i}= (\textbf{J}_k\textbf{W}_k\Sigma_i\textbf{W}_k^\top\textbf{J}_k^\top)[:2,:2]$. The $\textbf{W}_k$
is a viewing transformation matrix for input image $k$ and $\textbf{J}_k$ is the affine approximation of the projective transformation.
${\Sigma'_i}$ represents the covariance matrix in the 2D ray space.

The depth $\tilde{D}$ and normal $\tilde{N}$ for each pixel can also be
calculated via Gaussian splatting and alpha-blending:
\begin{equation}
        \tilde{D}=\frac{1}{1-T_{n+1}}\sum_{i=0}^{n}T_i\alpha_id_i(\textbf{u}),
\end{equation}
\begin{equation}
        \tilde{N}=\frac{1}{1-T_{n+1}}\sum_{i=0}^{n}T_i\alpha_i\textbf{R}_i[:,2].
\label{eq:gs_normal}
\end{equation}

Specifically, according to~\gs, the depth of pixel $\textbf{u}$ for each Gaussian kernel $i$
is computed by calculating the intersection of the ray cast through pixel $\textbf{u}$ 
with the Gaussian ellipse during splatting. So $d_i(\textbf{u})$ can be represented by Taylor expansion as:
\begin{equation}
    d_i(\textbf{u}) = d_i(\textbf{u}_i) + (\textbf{W}_k\textbf{R}_i)[2,:]\textbf{J}_{pr}^{-1}(\textbf{u}-\textbf{u}_i),
\end{equation}
where $\textbf{J}_{pr}^{-1}$ is the Jacobian inverse mapping one pixel from image space to tangent plane of the Gaussian surfel
as in~\cite{zwicker2001surface}, and $(\textbf{W}_k\textbf{R}_i)$ transforms the rotation matrix of a Gaussian surfel
to the camera space.

\subsection{Polarimetric image formation model}
Polarization cameras can capture polarized images in a single shot, 
which can be represented as $4$ different polarized angles of images 
$I=[I_0,I_{45},I_{90},I_{135}]$. 
The Stokes vector $S = \{s_0,s_1,s_2,s_3\}$ can be calculated by:
\begin{equation}
    S = \begin{bmatrix}\frac{1}{2}
        (I_0+I_{45}+I_{90}+I_{135})\\
        I_{0}-I_{90}\\
        I_{45}-I_{135}\\
        0
    \end{bmatrix},
\end{equation}
where we assume the light source is no circularly polarized here,
thus the $s_3$ is 0. 


According to \pandora and \nersp, we assume the incident environmental illumination is unpolarized, 
the Stokes vector 
for the incident light direction $\bm{\omega}$ can be expressed as:
\begin{equation}
    \V{s}_i(\bm{\omega}) = L(\bm{\omega}) [1, 0, 0, 0]^\top,
    \label{eq:env_stok}
\end{equation}
where $L(\bm{\omega})$ represents the light intensity. The polarization 
camera captures outgoing light that becomes partially polarized due to 
reflection, which is modeled using a $4 \times 4$ Muller matrix $\V{H}$. 
The outgoing Stokes vector $\V{s}$ then is formulated as the integral 
of the incident Stokes vector multiplied by this Mueller matrix:
\begin{equation}
    \V{s}(\V{v}) = \int_{\Omega} \V{H} \V{s}_i(\bm{\omega}) \,\, d\bm{\omega},
    \label{eq:out_stok}
\end{equation}
where $\V{v}$ indicates the view direction and $\Omega$ is the integral domain. 
Following the polarized BRDF (pBRDF) model~\cite{baek2018simultaneous}, the 
output Stokes vector can be divided into diffuse and specular components, 
represented by $\V{H}_d$ and $\V{H}_s$ respectively. $\V{s}$ 
can be represented as:
\begin{equation}
    \V{s}(\V{v}) = \int_{\Omega} \V{H}_d \V{s}_i(\bm{\omega}) \,\, d\bm{\omega} + \int_{\Omega} \V{H}_s \V{s}_i(\bm{\omega}) \,\, d\bm{\omega}.
    \label{eq:decomp_stok}
\end{equation}

Based on the derivations from \pandora and \nersp, the output Stokes vector can be 
further specified as:
\begin{eqnarray}
    \V{s}(\V{v}) =   L_d \begin{bmatrix}
        T_o^+ \\
        T_o^- \cos(2 \phi_n) \\
        - T_o^- \sin(2 \phi_n) \\
        0
    \end{bmatrix} +  L_s \begin{bmatrix}
        R^+ \\
        R^- \cos(2 \phi_h) \\
        -R^- \sin(2 \phi_h) \\
        0
    \end{bmatrix},
    \label{eq:photometric_cue}
\end{eqnarray}
where $L_d = \int_{\Omega} \rho L(\bm{\omega})  \bm{\omega}^\top \V{n} \, T_i^+T_i^-   \, d\bm{\omega}$ 
denotes the diffuse radiance associated with the surface normal $\V{n}$, 
Fresnel transmission coefficients~\cite{baek2018simultaneous} $T_{i, o}^+$ and 
$T_{i, o}^-$. The diffuse albedo is represented by $\rho$, and $\phi_n$
is the azimuth angle of the incident light. 

Similarly, $L_s = \int_{\Omega} L(\bm{\omega})  \frac{DG}{4 \V{n}^\top\V{v}}   \,\, d\bm{\omega}$ 
denotes the specular radiance, which involves Fresnel reflection coefficients~\cite{baek2018simultaneous} 
$R^+$ and $R^-$, and the incident azimuth angle $\phi_h$ concerning the half vector 
$\V{h} = \frac{\bm{\omega} + \V{v}}{\|\bm{\omega} + \V{v}\|_2^2}$. For the sake of simplification in our model, we assume that $\phi_h$
is equivalent to $\phi_n$. The Microfacet model incorporates the normal distribution and shadowing terms represented by 
$D$ and $G$~\cite{walter2007microfacet}. 



\section{Proposed method}
Our PolGS framework integrates multi-view polarized images, their corresponding masks, and camera pose information to produce a rich output that includes diffuse and specular components represented as Stokes vectors across various views, a reconstructed geometric mesh, and estimated environment light. In this section, we will analyze the surface normal representations employed in SDF-based approaches and 3DGS methods. Then we introduce a theoretical foundation for our polarization-guided reconstruction pipeline.
\vspace{-0.7em}
\begin{figure}
	\centering
	\includegraphics[width=\linewidth]{image/analysis.pdf}
	
	\caption{Comparison between SDF-based and 3DGS-based method on geometry representation.  (\textbf{Top}) The surface normal of a point has a strong relationship with its opacity in NeuS~\cite{wang2021neus}. 
		(\textbf{Bottom}) The surface normal of a point is dependent on its opacity in \gs.}
	\label{fig:normal-analysis}
    \vspace{-0.5em}
\end{figure}
\subsection{Analysis of surface normal representation}
In implicit neural networks, they use signed distance to determine if points are on the surface 
and the surface normals are obtained by calculating gradients of the SDF. However, 3DGS methods often treat the surface normal as inherent properties typically.
In this section, we analyze the differences
between the surface normal representations in SDF-based methods and 3DGS, explaining 
the reason we use polarimetric cues in the Gaussian Splatting method.
\vspace{-0em}
\begin{figure*}
	\centering
	\includegraphics[width=\linewidth]{image/pipeline.pdf}
	\caption{Pipeline of PolGS. We re-rendered Stokes vectors $\hat{s}$ by using the diffuse color $C$ from 3DGS and specular color $L_r$ from Cubemap encoder module, which is supervised by the ground truth Stokes information. }
	\label{fig:pipeline}
	\vspace{-1em}
\end{figure*}
\paragraph{SDF-based representation}
The implicit neural network approach, exemplified by NeuS~\cite{wang2021neus}, enhances the constraints on the reconstructed surface by incorporating a Signed Distance Function (SDF). As illustrated at the top of \fref{fig:normal-analysis}, the SDF value \( f(\textbf{x}) \), generated by the SDF Network, can be used to compute the normal vector \( \textbf{n}(\textbf{x}) \) at point \( \textbf{x} \). According to the definition, the opacity \( \rho(\textbf{x}) \) of a point \( \textbf{x} \)  is given by:

\begin{equation}
    \rho(\textbf{x}) = \max\left(\frac{-\frac{d\Phi_s}{dx}(f(\textbf{x}))}{\Phi_s(f(\textbf{x}))}, 0\right),
\label{eq:neus_density}
\end{equation}
where \( \Phi_s \) is a sigmoid function.

During optimization, the surface normal and the opacity at a given point influence each other. Consequently, after optimization, points farther from the surface tend to have lower opacity values, while the final normal vector is predominantly determined by points located on the surface. This mutual interaction results in a more precise and accurate representation of the surface.


\vspace{-1em}
\paragraph{3DGS-based representation}
In the case of \gs, surface normal is derived from the vector perpendicular to the plane of a 2D ellipsoid. As shown at the bottom of \fref{fig:normal-analysis}, for a single Gaussian kernel, there is no inherent relationship between opacity and the surface normal. 
The pixel-level constraint in 3DGS creates a probabilistic representation where no single Gaussian kernel is definitively assigned to a surface. Consequently, multiple Gaussian kernel configurations can potentially represent equivalent surface geometries, introducing inherent reconstruction ambiguity. 
The optimization mechanism effectively blends contributions from multiple Gaussian points, which offers reconstruction adaptability but compromises the precision of individual point geometric characterizations.

It is obvious that using the normal prediction model as a prior constraint in Gaussian Splatting benefits object reconstruction, as demonstrated by \gs. However, the normal prediction model is not always reliable especially in reflective cases, highlighting the need for a more robust method to provide prior information.
\vspace{-1em}
\paragraph{Polarimetric information for reflective surfaces}
Reflective surface reconstruction is challenging due to the view-dependent appearance. Unlike traditional RGB image inputs, 
polarimetric information can effectively constrain the surface normal during the rendering of Stokes Vectors with the pBDRF model~\cite{baek2018simultaneous}.
 Specifically, as demonstrated in \eref{eq:photometric_cue}, the diffuse and specular components of $s_1$ and $s_2$ 
exhibit a strong correlation with the object's surface normals. Leveraging this property, we incorporate polarimetric information as a prior in shape reconstruction and use the 3DGS-based method to accelerate this process.




\subsection{PolGS}
\paragraph{Network structure}
In this section, we introduce our proposed PolGS, which is a novel 3D reconstruction method that integrates
the Gaussian Splatting method with polarimetric information. The network structure of PolGS is shown in \fref{fig:pipeline}.
The framework comprises two primary components: the Gaussian Surfels module and the Cubemap Encoder module. 
Initially, the Gaussian Surfels module is employed to estimate the diffuse component of the object, 
seems like \gs. Subsequently, we utilize the CubeMap Encoder 
to assess the specular component, akin to the approach taken in \dr. While the CubeMap Encoder does not provide the roughness component, it effectively handles reflective or rough surfaces and maintains high computational efficiency due to its CUDA-based implementation. 
To enhance the rendering process, we incorporate a pBRDF model into the rendering formulation. 
This addition introduces a polarimetric constraint that further refines the Gaussian Splatting method, 
enabling more accurate and realistic 3D reconstructions.
The final rendering formulation model following 
\eref{eq:photometric_cue}
can be represented as:
\begin{equation}
    \hat{S} = C \begin{bmatrix}
        T_o^+ \\
        T_o^- \cos(2 \phi_n) \\
        - T_o^- \sin(2 \phi_n) \\
        0
    \end{bmatrix}+ L_r\begin{bmatrix}
        R^+ \\
        R^- \cos(2 \phi_h) \\
        -R^- \sin(2 \phi_h) \\
        0
    \end{bmatrix},
\label{eq:final_stokes}
\end{equation}
where $C$ is the diffuse color after Gaussian Surfels rendering according to \eref{eq:gs_render} and 
$L_r$ is the specular color after CubeMap rendering $E(\cdot)$, which can be represented as 
$L_r = E(2(\V{v}\cdot \textbf{n})\textbf{n}-\V{v})$.
\vspace{-1em}
\paragraph{Adjustment of spherical harmonic coefficients}
Due to \eref{eq:final_stokes}, the diffuse color rendered by Gaussian Surfels must remain consistent across different viewing directions. To satisfy this constraint, we adjust the spherical harmonic (SH) coefficients of the Gaussian Surfels to zero-order, ensuring the splatting results are view-independent.

\subsection{Training}
The overall training loss in PolGS formulated as a comprehensive weighted sum of multiple loss components:
\begin{equation}
	\mathcal{L} = \mathcal{L}_{rgb}+\lambda_{1}\mathcal{L}_{pol}+\lambda_{2}\mathcal{L}_{m}+\lambda_{3}\mathcal{L}_{o}	+\lambda_{4}\mathcal{L}_{d},
\end{equation}
where we set $\lambda_{1}=1$, $\lambda_{2}=0.1$, $\lambda_{3}=0.01$, 
$\lambda_{4}=0.01+0.1\cdot (iteration/15000)$ to balance the loss function.
\vspace{-1em}
\paragraph{Rendering Stokes loss $\mathcal{L}_{rgb}$ and $\mathcal{L}_{pol}$}
The rendering Stokes loss is combined with the $s_0$ (unpolarized image) rendering loss as 3DGS~\cite{kerbl20233d} and the 
$s_1$, $s_2$ rendering loss as \pandora~\cite{baek2018simultaneous}. These two loss functions can be represented as:
\begin{align}
        &\mathcal{L}_{rgb} = 0.8\cdot L_{1} (s_0,\hat{s}_0)+0.2\cdot L_{DSSIM}(s_0,\hat{s}_0),\\
        &\mathcal{L}_{pol} =L_{1} (s_1,\hat{s}_1)+L_{1} (s_2,\hat{s}_2).
\end{align}
\textbf{Mask loss $\mathcal{L}_{m}$}\quad
The mask loss is used to make the rendering results of the object more accurate, which can be represented as:
\begin{equation}
    \mathcal{L}_{m} = \Sigma\text{BCE}(\textbf{M},\hat{\textbf{M}}).
\end{equation}
\textbf{Opacity loss $\mathcal{L}_{o}$}\quad
The opacity loss follows \gs to encourage the opacity of the Gaussian points to be close to 1 or 0. It can be represented as:
\begin{equation}
    \mathcal{L}_{o} = \Sigma\text{exp}(-{20(o_i-0.5)}^2).
\end{equation}
\textbf{Depth-normal consistency loss $\mathcal{L}_{d}$}\quad
The depth-normal consistency loss follows \gs to make the rendered depth and normal of the object to be more consistent.
It can be represented as:
\begin{equation}
    \mathcal{L}_{d} = 1-\hat{\textbf{N}}\cdot N(V(\hat{\textbf{D}})),
\end{equation}
where $V(\cdot)$ transforms each pixel and its depth to a 3D point and $N(\cdot)$ calculates the normal from neighboring points using the cross product.


